\chapter{Solvers}
\label{ch:solvers}

All solvers read the same \wtwof{} input files (\file{constant/waveProperties}, relaxation
zones, probes) and accept every wave theory, including the HOS forcing of Chapter~\ref{ch:hos}.
They differ in the interface-capturing scheme and in how a floating body is moved.

\begin{table}[htbp]
\centering\small
\begin{tabularx}{\linewidth}{@{}l l l X@{}}
\toprule
Solver & Interface & Wave zones & Body motion \\
\midrule
\app{waveFoam}                & MULES       & yes & \kw{sixDoFRigidBodyMotion} \\
\app{waveIsoFoam}             & isoAdvector & yes & \kw{sixDoFRigidBodyMotion} \\
\app{floatStepperWaveIsoFoam} & isoAdvector & yes & FloatStepper \\
\app{floatStepper}            & isoAdvector & no  & FloatStepper \\
\bottomrule
\end{tabularx}
\caption{Solvers in \fwt{}.}
\label{tab:solvers}
\end{table}

\section{\texorpdfstring{\app{waveFoam}}{waveFoam}}
\app{waveFoam} is the \wtwof{} solver (\wtwof{} manual, Sec.~5.1), ported to OpenFOAM v2506
and with the \kw{ddtCorr} switch of Sec.~\ref{sec:ddtcorr}. It supports dynamic meshes, so a
rigid body can be moved with OpenFOAM's \kw{sixDoFRigidBodyMotion} solver, with or without
MoorDyn moorings (Sec.~\ref{sec:moordyn-sixdof}).

\section{\texorpdfstring{\app{waveIsoFoam}}{waveIsoFoam}}
\label{sec:waveisofoam}
\app{waveIsoFoam} is \app{waveFoam} with the volume fraction advected by the geometric
isoAdvector scheme (as in OpenFOAM's \app{interIsoFoam}) instead of MULES. isoAdvector keeps the
interface sharp over long wave propagation distances and is the scheme FloatStepper is built
on. The relaxation zones are applied after the advection step, as in \app{waveFoam}.

The case set-up is that of \app{waveFoam}, except for the \kw{alpha.water} entry in
\file{system/fvSolution}:
\begin{lstlisting}[caption={isoAdvector settings in \file{system/fvSolution} (\file{tutorials/waveIsoFoam}).}]
"alpha.water.*"
{
    reconstructionScheme plicRDF;   // isoAlpha (faster) or plicRDF (more accurate)
    isoFaceTol      1e-6;
    surfCellTol     1e-3;
    nAlphaBounds    3;
    snapTol         1e-12;
    clip            true;
    nAlphaSubCycles 1;              // required
    cAlpha          1;              // not used by isoAdvector
    // MULES entries (nAlphaCorr, MULESCorr, nLimiterIter) are ignored
}
\end{lstlisting}
The meaning of these entries is that of OpenFOAM's \app{interIsoFoam}. Keep the Courant number
below about 0.5 (\kw{maxCo} and \kw{maxAlphaCo} in \file{system/controlDict}); isoAdvector is
not stable at large Courant numbers.

\section{The \texorpdfstring{\kw{ddtCorr}}{ddtCorr} switch}
\label{sec:ddtcorr}
In the pressure equation, OpenFOAM adds the term \kw{fvc::ddtCorr(U, phi)} to the predicted face
flux. This Rhie--Chow-type correction for transient flows uses the difference between the face
flux and the interpolated cell velocity of the previous time step, keeping the two consistent.
It changes the flux near the free surface and near moving bodies, so it can be useful to test a
case with and without it. In \fwt{}, all solvers of
Table~\ref{tab:solvers} let this term be switched off from \file{system/fvSchemes}:
\begin{lstlisting}[caption={The \kw{ddtCorr} switch in \file{system/fvSchemes}.}]
ddtCorr  on;    // on (default) | off

ddtSchemes
{
    default Euler;
}
\end{lstlisting}
The entry is at the top level of \file{fvSchemes}, not inside \kw{ddtSchemes}. Without the entry,
the term is included, which is the standard OpenFOAM behaviour. The switch is read in every
pressure corrector, so it can be changed while the case is running.

\section{Floating bodies with FloatStepper}
\label{sec:floatstepper}
\subsection{Method}
FloatStepper~\cite{roenby2024robust} removes the added-mass instability of light bodies in heavy
fluids. The force (and torque) on the body is split into an added-mass part and the rest,
\begin{equation}
  \mathbf F = \mathbf F_\mathrm{other} - \mathsf A\,\mathbf a
  \quad\Rightarrow\quad
  \mathbf a = (\mathsf M + \mathsf A)^{-1}\,\mathbf F_\mathrm{other},
\end{equation}
where $\mathsf M$ is the body mass matrix, $\mathsf A$ the added-mass matrix and $\mathbf a$
the body acceleration. $\mathbf F_\mathrm{other}$ is found from a test step with zero body
acceleration, and $\mathsf A$ is computed numerically as the force per unit acceleration. No
outer correctors or acceleration relaxation are needed.

\fwt{} contains FloatStepper as a git submodule in
\file{applications/solvers/floatStepperWaveFoam} and builds two solvers from it:
\begin{description}
  \item[\app{floatStepperWaveIsoFoam}] FloatStepper with the \wtwof{} relaxation zones and wave
    theories, including HOS. This is the solver for floating bodies in waves.
  \item[\app{floatStepper}] the original FloatStepper solver, without relaxation zones.
\end{description}
Both support the \kw{ddtCorr} switch (Sec.~\ref{sec:ddtcorr}).

\subsection{Case set-up}
A FloatStepper case differs from a \kw{sixDoFRigidBodyMotion} case in four places.

\paragraph{1. Solver and libraries} in \file{system/controlDict}:
\begin{lstlisting}
application     floatStepperWaveIsoFoam;
libs            (moorDynForFloatStepper waveModels);  // MoorDyn coupling, if used
\end{lstlisting}

\paragraph{2. Mesh motion} in \file{constant/dynamicMeshDict}:
\begin{lstlisting}[caption={FloatStepper mesh motion (\file{tutorials/floatStepperWaveIsoFoam}).},label={lst:fsdyn}]
dynamicFvMesh   dynamicPointDisplacementFvMesh;

solvers
{
    box                                 // body name, as in 0/uniform/floaters
    {
        motionSolverLibs ( floaterMotion );
        motionSolver     floaterMeshMotionSolver;
        floaterMeshMotionSolverCoeffs
        {
            patches         (box);
            innerDistance   0.075;      // rigid mesh motion inside
            outerDistance   0.4;        // no mesh motion outside
            linDirs         (1 0 1);    // translation in x and z only
            rotDirs         (0 1 0);    // rotation about y only
            addedMassModel  pisoAddedMass;
            MaddUpdateFreq  1;          // added mass every n-th time step
            restraints
            {
                // e.g. MoorDyn, Chapter 5
            }
        }
    }
}
\end{lstlisting}
\kw{linDirs} and \kw{rotDirs} replace the \kw{constraints} of \kw{sixDoFRigidBodyMotion}; by
default all six degrees of freedom are free. Computing the added mass is expensive;
$\kw{MaddUpdateFreq}=n$ recomputes it only every $n$-th time step. The entries
\kw{constraints}, \kw{solver}, \kw{accelerationRelaxation} and \kw{accelerationDamping} of
\kw{sixDoFRigidBodyMotion} are not used. Restraints have the same syntax as in
\kw{sixDoFRigidBodyMotion}, with \kw{sixDoFRigidBody} replaced by \kw{floater} in the class
names (e.g.\ \kw{floaterMotionRestraint}).

\paragraph{3. Body properties} in \file{0.orig/uniform/floaters}, in a dictionary with the body
name used in \file{dynamicMeshDict}:
\begin{lstlisting}[caption={Body properties, \file{0.orig/uniform/floaters}.}]
box
{
    centreOfRotation    (0 0 -0.0126);
    orientation         (1 0 0 0 1 0 0 0 1);
    velocity            (0 0 0);
    acceleration        (0 0 0);
    angularVelocity     (0 0 0);
    angularAcceleration (0 0 0);
    mass                3.148;
    centreOfMass        $centreOfRotation;
    momentOfInertia     (0.015 0 0 0.0167 0 0.021);   // xx xy xz yy yz zz
    initialCentreOfRotation $centreOfRotation;
    initialOrientation      $orientation;
    initialCentreOfMass     $initialCentreOfRotation;
}
\end{lstlisting}
Mass properties given in \file{dynamicMeshDict} are not used.

\paragraph{4. Velocity on the body} in \file{0.orig/U}:
\begin{lstlisting}
box
{
    type    floaterVelocity;
    slip    true;            // default: no-slip
    value   uniform (0 0 0);
}
\end{lstlisting}

\subsection{Output}
The body motion is written by the \kw{floaterStateWriter} function object:
\begin{lstlisting}[caption={Writing the body motion (\file{system/controlDict}).}]
functions
{
    floaterMotion
    {
        type            floaterStateWriter;
        libs            (floaterStateWriter);
        writeControl    timeStep;
        writeInterval   1;
        bodyName        box;
    }
}
\end{lstlisting}

\Warn{FloatStepper supports one rigid body, a morphing mesh (no overset) and the isoAdvector
interface scheme. See the FloatStepper README for further details and limitations.}
